\documentclass[../main.tex]{subfiles}
\begin{document}

\section{Problem Statement}
\label{sec:problem}

Genome-based prediction of antimicrobial resistance (\gls{AMR}) in \textit{Salmonella} is the subject of this research, which builds on and critically re-evaluates the expert-marker prediction study of the original 2025 work \cite{baseline2025}. \gls{AMR} is one of the greatest threats to global public health, estimated to be associated with millions of deaths each year \cite{murray2022global}. \textit{Salmonella} is a common food-borne pathogen and a suitable subject for studying AMR because large volumes of genomic data with paired resistance phenotypes are available. The problem addressed here is: given the genomic data of a \textit{Salmonella} strain, predict whether that strain is resistant or susceptible to a specific antibiotic.

Accurate and reliable genome-based AMR prediction can support rapid diagnosis, epidemiological surveillance and early warning, reducing reliance on time-consuming conventional susceptibility testing. However, the problem has three characteristic difficulties: the resistant class is often rare (some antibiotics have only about 6\% resistant samples), the genomic feature space is very large (tens of thousands of genes), and the resistance mechanisms differ across antibiotics.

\section{Background and Problems of Research}
\label{sec:background}

Many studies have applied machine learning to predict AMR from \textit{Salmonella} genomes, using gene-based features (the accessory genome), point variants (core \gls{SNP}), and \gls{MIC} values \cite{nguyen2019using}. Most report high accuracy on random data splits.

Nevertheless, current solutions share three important limitations. First, with strongly imbalanced data, reporting accuracy alone is misleading and is rarely accompanied by statistical significance testing. Second, when many feature/model configurations are tried and the best one is chosen by its score on the test set, the results are inflated by a ``winner's curse'' effect. Third, bacterial genomes have strong population structure: strains of the same phylogenetic lineage are highly similar, so a random split may let the model learn lineage identity rather than the true resistance mechanism, giving an optimistic estimate of generalization \cite{earle2016identifying}.

\section{Research Objectives and Conceptual Framework}
\label{sec:objectives}

The objective of this research is to build and \textit{reliably evaluate} an AMR prediction framework for \textit{Salmonella}. The research follows an antibiotic-aware adaptive feature fusion approach: instead of a fixed pipeline, each antibiotic is assigned a feature module and model suited to its resistance mechanism.

Corresponding to the three limitations in Section~\ref{sec:background}, the research proposes: (i) rigorous statistical testing (a paired t-test on repeated cross-validation folds, with a Wilcoxon signed-rank test and a bootstrap confidence interval) to determine whether the advantage of adaptive fusion is real or noise; (ii) evaluation along true phylogenetic lineages (leave-lineage-out based on internal core SNPs and external \gls{MLST}) to detect confounding; and (iii) true mechanism-level mutation calling (\gls{QRDR}) for the quinolone group to compare against annotation features.

\section{Contributions}

This research makes five main contributions:

\begin{enumerate}
\item It confirms that the accessory genome is the dominant AMR predictive signal, consistently outperforming core SNPs across the five antibiotics.
\item It proposes an antibiotic-aware adaptive feature fusion method that improves F1 on 4 of 5 drugs over the strong expert-marker baseline, with each antibiotic selecting a different best module, and validates this with repeated cross-validation and a paired t-test.
\item It adds lineage-aware evaluation through true phylogenetic splits, identifying which antibiotics carry a robust mechanism-driven signal and which are more lineage-dependent.
\item It performs true QRDR mutation calling (gyrA, parC) for the quinolone group and shows that mechanism features generalize across lineages better than annotation features.
\item It evaluates the deployability of the model (probability calibration, screening utility) and provides reproducible notebooks together with a consolidated report. All source code, notebooks and the report are publicly available\footnote{\url{https://github.com/bennyy117/adaptive_salmonella}}.
\end{enumerate}

\section{Organization of Research}

The remainder of the research is organized as follows. Chapter~\ref{ch:lit} analyses the problem context and related work and presents the biological and machine-learning background. Chapter~\ref{ch:method} describes the methodology. Chapter~\ref{ch:arch} presents the proposed architecture and pipeline with the input/output matrix sizes and data formats. Chapter~\ref{ch:impl} explains the source code of each pipeline block line by line. Chapter~\ref{ch:theory} gives the theoretical analysis of the paired t-test and formalizes population-structure confounding. Chapter~\ref{ch:results} presents and discusses the experimental results, organized by research question. Chapter~\ref{ch:concl} summarizes the results and limitations and suggests future work.

\end{document}
